\chapter{The First Words: Variables, Memory Boxes, and Scope}

\section{What is a Variable?}

\begin{kidbox}{The Labeled Cardboard Box}
Imagine you have a messy room full of toys. You want to save your shiny red race car so you don't lose it.
What do you do? You grab an empty cardboard box.
With a big black crayon, you write a name on the outside: \textbf{"my\_car"}.
Then, you put the red car inside the box and place it on your shelf!

Whenever you want to play with the car, you don't have to search the whole house. You just look at the shelf, find the box labeled \textbf{"my\_car"}, open it up, and there it is!
In Shae, we write:
\begin{lstlisting}[language=shaelang]
let my_car = "Red Race Car";
\end{lstlisting}
\end{kidbox}

\section{Immutability: Open Boxes vs. Superglued Boxes}

\begin{kidbox}{The Superglue Trick}
Now imagine you built a precious LEGO spaceship that took three weeks to make. You put it in a box labeled \textbf{"spaceship"}.
You tell your little brother: "Do NOT touch this! Do NOT swap it with a potato!"
To be completely safe, you put superglue all around the lid. Once closed, nobody can EVER replace the toy inside!
In Shae, by default, our boxes are safe and sealed! If you try to swap out the toy without permission, Shae stops you:
\begin{lstlisting}[language=shaelang]
let score = 10;
// score = 20; // Error: Cannot reassign immutable variable!
\end{lstlisting}
\end{kidbox}

\begin{deepbox}{Stack Allocation, Registers, and SSA Form}
At the hardware level, variables do not have human names like \texttt{"my\_car"} or \texttt{"score"}. The CPU operates purely on:
\begin{enumerate}
    \item \textbf{General Purpose Registers}: High-speed 64-bit storage locations directly on the CPU die (e.g., \texttt{\%rax}, \texttt{\%rcx}, \texttt{\%rdi} on x86\_64; \texttt{x0} through \texttt{x30} on ARM64). Access latency is sub-nanosecond ($<0.5\text{ ns}$).
    \item \textbf{The Call Stack}: Contiguous virtual memory managed by the Stack Pointer (\texttt{\%rsp}). Local variables are assigned fixed offsets from the Base Pointer (\texttt{\%rbp}). Allocation is a single instruction (\texttt{sub rsp, N}).
    \item \textbf{The Heap}: Dynamic memory managed by an allocator (\texttt{malloc}, \texttt{jemalloc}, or system virtual memory pages via \texttt{mmap}).
\end{enumerate}
In the Shae Bytecode VM, local variables are resolved at compile time to integer stack slot offsets (\texttt{OpGetLocal(slot)}, \texttt{OpSetLocal(slot)}). This eliminates runtime string hash lookups entirely!
\end{deepbox}

\section{Destructuring Bindings: Unpacking Surprise Eggs}

\begin{kidbox}{Cracking Open the Surprise Egg}
Have you ever opened a chocolate surprise egg that had three little toys inside?
Instead of pulling out the whole egg and carrying it around, you crack it open on the table:
"Here is the sticker, here is the spinning top, and here is the mini dinosaur!"
In Shae, you can unpack arrays and maps in one clean move:
\begin{lstlisting}[language=shaelang]
// Unpacking an array of favorite colors
let [first, second, ..rest] = ["Blue", "Green", "Yellow", "Purple"];

println("First favorite: " + first);   // Blue
println("Second favorite: " + second); // Green
println("The others: " + rest);       // ["Yellow", "Purple"]
\end{lstlisting}
\end{kidbox}

\begin{syntaxbox}{Destructuring Patterns in Let Statements}
Shae supports rich destructuring patterns across both arrays and key-value maps:
\begin{lstlisting}[language=shaelang]
// 1. Array destructuring with rest collector
let [head, ..tail] = [1, 2, 3, 4, 5];

// 2. Map / Struct destructuring with rest collector
let user = {
    "name": "Shae Dev",
    "role": "Engineer",
    "level": 42,
    "active": true
};
let { name, role, ..extra } = user;

println("Name: " + name);   // Shae Dev
println("Role: " + role);   // Engineer
println("Extra: " + extra); // {"level": 42, "active": true}
\end{lstlisting}
The compiler validates at compile-time that rest identifiers do not collide with bound symbols.
\end{syntaxbox}

\section{Scope and Lifetime: The Playroom Rules}

\begin{kidbox}{What Happens in the Playroom Stays in the Playroom}
When you walk into your playroom, you bring out your crayons and draw pictures on the playroom table.
When mom says: "Dinner time! Leave the playroom!", you step out into the hallway.
Can you still touch the crayons that are inside the playroom? No! The door is closed!
And while you are eating dinner, the cleanup robot tidies up the room.
In Shae, every set of curly braces \texttt{\{ \}} is like a little room (we call it a \textbf{Scope}):
\begin{lstlisting}[language=shaelang]
let outside = "Sun is shining";

{
    let secret_toy = "Golden Key";
    println(secret_toy); // Works! We are inside the room!
    println(outside);    // Works! You can look out the window!
}

// println(secret_toy); // ERROR! The room was cleaned up!
\end{lstlisting}
\end{kidbox}
